\documentclass{article}
\usepackage[T1]{fontenc}
\usepackage{lmodern}
\usepackage{graphicx}
\usepackage{amsmath,amssymb}
\usepackage[a4paper,margin=2cm]{geometry}
\usepackage[absolute,overlay]{textpos}
\setlength{\TPHorizModule}{1mm}
\setlength{\TPVertModule}{1mm}
\usepackage[indonesian]{babel}
\usepackage{hyperref}
\setlength{\emergencystretch}{2em}
% unit: Halaman 14

\begin{document}

% === Halaman 14 (python/native) ===

• RC5 dibuat oleh Ronald Rivest dari Laboratorium RSA  pada tahun 1994.

• Salah satu finalis AES, yaitu RC6, dikembangkan dari RC5 

• Tidak seperti algoritma cipher blok lainnya, RC5 mempunyai:

 - ukuran blok yang variabel (16, 32, 64)

 - panjang kunci yang variabel (0 sampai 2040 bit)

 - dan jumlah putaran yang variabel (0 sampai 255). 

14


Parameter 

Simbol 

Nilai yang dibolehkan 

Ukuran blok (dalam bit) 

w 

16, 32, 64 

Jumlah putaran 

r 

0, 1, …, 255 

Panjang kunci eksternal K 

(dalam byte, 1 byte = 8 bit) 

b 

0, 1, …, 255 


RC5

\end{document}
